% Generated draft: requires mathematical and source review.
\documentclass[11pt]{article}
\usepackage[margin=25mm]{geometry}
\usepackage{fontspec}
\setmainfont{Latin Modern Roman}
\usepackage{amsmath,amssymb,mathtools,unicode-math}
\setmathfont{Latin Modern Math}
\usepackage{enumitem,needspace,xurl,hyperref}
\setlength{\parindent}{0pt}
\setlength{\parskip}{5pt}
\setlength{\emergencystretch}{3em}
\newcommand{\N}{\mathbb{N}}
\newcommand{\Z}{\mathbb{Z}}
\newcommand{\Q}{\mathbb{Q}}
\newcommand{\R}{\mathbb{R}}
\newcommand{\C}{\mathbb{C}}
\newcommand{\F}{\mathbb{F}}
\newcommand{\A}{\mathbb{A}}
\newcommand{\PP}{\mathbb{P}}
\newcommand{\E}{\mathbb{E}}
\newcommand{\eps}{\varepsilon}
\newcommand{\dd}{\,\mathrm{d}}
\begin{document}
\section*{UnsolvedMath ID: 800365}
\textbf{Source problem code:} catalogue entry
\subsection{Definitions and mathematical statement}
Let \((W_i,S_i)\), for \(i\in\{1,2\}\), be Coxeter systems, with length functions \(\ell_i\) and Bruhat orders \(\le_i\). Here Bruhat order is characterized by the reduced-subword criterion. No finiteness assumption on either Coxeter group is supplied. For \(J\subseteq S\), let \(W_J=\langle J\rangle\) be the standard parabolic subgroup. The notation \(W^J\) denotes distinguished minimal-length parabolic coset representatives, with the convention qualification below. For \(a,b\in W\), set
\[
[a,b]=\{x\in W:a\le x\le b\},\qquad [a,b]^J=[a,b]\cap W^J,
\]
with their induced Bruhat orders. An interval is empty when its endpoints are not comparable in the indicated direction.

Let \(q\) be an indeterminate. Write \(P_{a,b}^{J,q}(q)\) for the \(q\)-type parabolic Kazhdan--Lusztig polynomial associated with \((W,S,J)\); the superscript \(q\) specifies the parabolic parameter, while the final \(q\) is the polynomial variable. Use the customary zero extension when \(a\not\le b\).

\textbf{Conjecture (Brenti, 2024, Conjecture 3).} For every pair of Coxeter systems \((W_1,S_1)\), \((W_2,S_2)\), every \(J_1\subseteq S_1\), \(J_2\subseteq S_2\), every \(u,v\in W_1^{J_1}\), \(w,z\in W_2^{J_2}\), and every poset isomorphism
\[
f:[u,v]\longrightarrow[w,z]
\]
satisfying
\[
f\bigl([u,v]^{J_1}\bigr)=[w,z]^{J_2},
\]
one has the polynomial identity
\[
P_{u,v}^{J_1,q}(q)=P_{w,z}^{J_2,q}(q).
\]

\textbf{Unresolved conventions.} The supplied extract does not define \(W^J\) or the normalization of the parabolic Kazhdan--Lusztig polynomials. Minimal representatives for right cosets \(W/W_J\) satisfy \(\ell(xs)>\ell(x)\) for all \(s\in J\); minimal representatives for left cosets \(W_J\backslash W\) satisfy \(\ell(sx)>\ell(x)\) for all \(s\in J\). The statement retains the source notation without selecting between these conventions. The representative convention and polynomial normalization must be checked together against the original source. The zero extension above makes the statement meaningful without adding endpoint-comparability hypotheses that were not explicitly supplied.
\subsection{Short English statement}
Brenti's conjecture asserts that the \(q\)-type parabolic Kazhdan--Lusztig polynomial is determined by the Bruhat interval together with its distinguished subset of minimal parabolic coset representatives. An isomorphism of the full intervals must carry that subset onto the corresponding subset. The supplied registry labels this conjecture open; its notation conventions still require source review.
\subsection{Sources}
\begin{itemize}
\item Local UnsolvedMath registry, permanent ID \(800365\), exact source problem code \texttt{catalogue entry}, title ``Coxeter Groups and Unimodality --- Conjecture 3.'' The record attributes the conjecture to Brenti, with proposed year 2024, and supplies the status label \texttt{open}. This label is source data, not an independently verified literature assessment.
\item Brenti, \emph{Some Open Problems on Coxeter Groups and Unimodality} (2024), Conjecture 3. Supplied source URL: \url{https://arxiv.org/abs/2410.09897}. The record reports access on 2026-10-16 and extraction method \texttt{arxiv\_latex}; these are supplied provenance details, not a claim of independent source access.
\end{itemize}
\end{document}
